\documentclass[11pt,a4paper]{report}
\usepackage[T1]{fontenc}
\usepackage[utf8]{inputenc}
\usepackage[french]{babel}
\usepackage{amsmath,amssymb,amsthm,mathtools}
\usepackage{geometry}
\usepackage{booktabs}
\usepackage{array}
\usepackage{longtable}
\usepackage{graphicx}
\usepackage{xcolor}
\usepackage{hyperref}
\usepackage{listings}
\usepackage{float}
\usepackage{enumitem}
% NOTE: natbib supprimé car on utilise thebibliography inline
\geometry{margin=2.5cm}

\hypersetup{
  colorlinks=true,
  linkcolor=blue,
  urlcolor=blue,
  citecolor=blue
}

\lstdefinestyle{matlabstyle}{
  language=Matlab,
  basicstyle=\ttfamily\small,
  keywordstyle=\color{blue!70!black},
  commentstyle=\color{green!40!black},
  stringstyle=\color{orange!60!black},
  numbers=left,
  numberstyle=\tiny\color{gray},
  stepnumber=1,
  frame=single,
  breaklines=true,
  tabsize=2,
  showstringspaces=false
}

\newtheorem{theorem}{Theoreme}[chapter]
\newtheorem{proposition}[theorem]{Proposition}
\newtheorem{lemma}[theorem]{Lemme}

\title{Analyse et simulation numerique de l'equation de Nagumo\\
\large Etude theorique, schemas numeriques et stabilite}
\author{Mawounan Jolemon GLODJO}
\date{22 avril 2026}

\begin{document}
\maketitle

\begin{abstract}
Ce rapport presente une etude complete de l'equation de Nagumo en dimension un,
modele reaction-diffusion bistable apparaissant en neurobiologie, en ecologie et dans
l'etude des fronts de propagation.
Nous commencons par la modelisation mathematique, l'analyse des equilibres, l'existence
de fronts voyageurs et la vitesse analytique associee.
Nous derivons ensuite plusieurs methodes de discretisation en temps et en espace:
Runge-Kutta d'ordre 4, Euler implicite, Adams-Bashforth d'ordre 2 et Adams-Moulton d'ordre 2.
L'analyse de stabilite est menee au niveau lineaire puis par une etude de Von Neumann,
afin de mettre en evidence les contraintes CFL et les differences entre schemas explicites
et implicites.
Enfin, une base de simulation MATLAB est fournie pour comparer les schemas, suivre la
propagation du front, verifier la convergence et visualiser les resultats.
\end{abstract}

\tableofcontents

\chapter*{Preambule -- Mise en contexte}
\addcontentsline{toc}{chapter}{Preambule -- Mise en contexte}
L'equation de Nagumo a ete introduite \cite{Nagumo1962} comme modele simple mais riche pour decrire la
propagation d'un signal dans un milieu bistable. Elle intervient dans des problemes de
migration cellulaire \cite{MurrayMathematicalBiology}, d'activation neuronale, de dynamique des populations et plus
largement dans les phenomenes de front de reaction-diffusion \cite{Volpert2000}.

L'interet principal de ce modele est qu'il combine trois ingredients fondamentaux:
la diffusion, qui regularise l'espace; la reaction locale, qui cree plusieurs etats
stables; et l'effet de seuil, qui determine la direction et la vitesse de propagation.

Le present rapport a pour objectif de construire une etude mathematique propre et une
chaine numerique complete. L'accent est mis sur la derivation des schemas, leur domaine
de stabilite, et leur comportement en simulation.

\chapter{Partie I -- Analyse theorique}
\section{Formulation mathematique}
On considere l'equation de Nagumo sur un domaine unidimensionnel $\Omega=[0,L]$:
\begin{equation}
\frac{\partial u}{\partial t} = D\,\Delta u + u(1-u)(u-a), \qquad x \in \Omega,\; t>0,
\end{equation}
ou, en dimension un,
\begin{equation}
\frac{\partial u}{\partial t} = D\,\frac{\partial^2 u}{\partial x^2} + u(1-u)(u-a).
\end{equation}
Ici $u(x,t)$ represente un etat d'activation, $D>0$ est le coefficient de diffusion
et $a \in (0,1)$ est un parametre de seuil.

On impose une condition initiale $u(x,0)=u_0(x)$ ainsi que des conditions aux bords
choisies selon le contexte. Dans ce rapport, on retient des conditions de Neumann
homogenes:
\begin{equation}
\partial_x u(0,t)=0, \qquad \partial_x u(L,t)=0,
\end{equation}
qui correspondent a un milieu isole sans flux de matiere aux extremites.

\section{Points d'equilibre et analyse qualitative}
On note
\begin{equation}
 f(u)=u(1-u)(u-a).
\end{equation}
Les points d'equilibre homogenes sont les racines de $f$, soit
\begin{equation}
 u^* \in \{0,a,1\}.
\end{equation}
Pour l'ODE associee sans diffusion, $\dot u=f(u)$, on a
\begin{equation}
 f'(u) = -3u^2 + 2(1+a)u - a.
\end{equation}
D'ou
\begin{equation}
 f'(0)=-a<0, \qquad f'(a)=a(1-a)>0, \qquad f'(1)=a-1<0.
\end{equation}
Ainsi $u=0$ et $u=1$ sont localement stables, tandis que $u=a$ est instable.

La fonction de reaction est bistable: elle pousse la solution vers l'un des deux etats
stables selon que l'etat initial se situe au-dessus ou en dessous du seuil.
Cette structure est typique des systemes a changement de phase ou a propagation de front.

\subsection*{Portrait de phase de l'EDO locale}
Sans diffusion, l'equation devient $\dot u=f(u)$.
Le signe de $f$ permet de lire le sens d'evolution:
\begin{itemize}[leftmargin=1.5em]
  \item si $0<u<a$, alors $f(u)<0$ et la solution decroit vers 0;
  \item si $a<u<1$, alors $f(u)>0$ et la solution cro\^{\i}t vers 1;
  \item au voisinage de $u=a$, la dynamique est repulsive.
\end{itemize}
Ce portrait de phase explique la coexistence de deux phases stables et d'un seuil de bascule.

\section{Solutions en fronts voyageurs : approche théorique}
\subsection*{Motivation et justification}
L'une des caractéristiques remarquables de l'équation de Nagumo en régime bistable
est l'existence de \textbf{solutions de type fronts voyageurs}, appelées aussi \textit{ondes progressives}.
Ces solutions représentent des transitions stables entre deux états d'équilibre qui se propagent
à vitesse constante dans l'espace.

\textbf{Pourquoi chercher de telles solutions?}
\begin{enumerate}[leftmargin=2em]
  \item Les fronts voyageurs sont des \textit{structures cohérentes} qui émergent naturellement
        des équations de réaction-diffusion bistables.
  \item Elles représentent des \textit{attracteurs globaux}: de nombreuses conditions initiales
        convergent vers ce type de comportement asymptotique.
  \item Elles sont \textit{observées expérimentalement} dans les systèmes biologiques
        (ondes nerveuses, fronts enzymatiques, etc.).
\end{enumerate}

\subsection*{Construction de l'équation pour les fronts voyageurs}
On cherche des solutions de la forme:
\begin{equation}
 u(x,t) = \phi(\xi), \qquad \xi = x - ct,
\end{equation}
où $c$ est la vitesse du front (à déterminer) et $\phi(\xi)$ est le profil spatial.
Cette approche transforme l'équation aux dérivées partielles en une équation différentielle ordinaire.

\textbf{Dérivation des relations de dérivation:}
En utilisant la règle de chaîne,
\begin{align}
 \frac{\partial u}{\partial t} &= \frac{d\phi}{d\xi} \cdot \frac{d\xi}{dt} = -c \phi'(\xi), \\
 \frac{\partial u}{\partial x} &= \phi'(\xi), \quad \frac{\partial^2 u}{\partial x^2} = \phi''(\xi).
\end{align}

\textbf{Substitution dans l'équation de Nagumo:}
En remplaçant ces dérivées dans
\begin{equation}
 \frac{\partial u}{\partial t} = D \frac{\partial^2 u}{\partial x^2} + f(u),
\end{equation}
on obtient:
\begin{equation}
 -c\phi'(\xi) = D\phi''(\xi) + \phi(\xi)(1-\phi(\xi))(\phi(\xi)-a),
\end{equation}

soit:
\begin{equation}\boxed{
 D\phi'' + c\phi' + f(\phi) = 0, \qquad \text{où } f(\phi)=\phi(1-\phi)(\phi-a).
}
\end{equation}

\subsection*{Conditions asymptotiques et connectivité}
Pour que le front connecte deux états stables de manière physiquement pertinente,
on impose les \textbf{conditions asymptotiques}:
\begin{equation}
 \phi(-\infty) = 1, \qquad \phi(+\infty) = 0.
\end{equation}

Ces conditions signifient que:
\begin{itemize}[leftmargin=1.5em]
  \item Loin derrière le front ($\xi \to -\infty$), la solution vaut 1 (état "élevé" stable).
  \item Loin devant le front ($\xi \to +\infty$), la solution vaut 0 (état "bas" stable).
  \item Le front assure une transition continue et monotone (ou oscillatoire amortie) entre ces deux états.
\end{itemize}

\textbf{Remarque sur les vitesses:} Toute vitesse $c$ n'est pas admissible. 
Il existe une \textbf{unique vitesse minimale} $c^*$ pour laquelle existe un profil $\phi$ vérifiant
les conditions asymptotiques et reste entre 0 et 1. C'est un résultat profond de la théorie
des équations de réaction-diffusion bistables (voir e.g., \cite{FifeMcLeod1977}).

\subsection*{Existence et unicité du profil}
\begin{theorem}[Existence et unicité du front]
Pour l'équation de Nagumo avec $0 < a < 1$ et $D > 0$, il existe une unique
vitesse critique
\begin{equation}
 c^* = \sqrt{2D}\left(\frac{1}{2} - a\right)
\end{equation}
et un unique profil (à translation) $\phi_*(\xi)$ tel que:
\begin{enumerate}[leftmargin=1.5em]
  \item $\phi_*$ satisfait $D\phi_*'' + c^*\phi_*' + f(\phi_*) = 0$;
  \item $\phi_*(-\infty)=1$, $\phi_*(+\infty)=0$;
  \item $0 < \phi_*(\xi) < 1$ pour tout $\xi \in \mathbb{R}$;
  \item $\phi_*$ décroît monotonement (profil régulier sans oscillations).
\end{enumerate}
\end{theorem}

\textbf{Observations:}
\begin{itemize}[leftmargin=1.5em]
  \item Si $a < \tfrac{1}{2}$, alors $c^* > 0$ : le front progresse vers la droite (de 1 vers 0).
  \item Si $a > \tfrac{1}{2}$, alors $c^* < 0$ : le front se propage vers la gauche.
  \item Si $a = \tfrac{1}{2}$, alors $c^* = 0$ : le front devient stationnaire.
  \item La vitesse dépend de $D$ et $a$ uniquement, pas de conditions initiales (propriété remarquable).
\end{itemize}

\subsection*{Forme explicite du profil}
Une approximation classique du profil de front est
\begin{equation}
 \phi(\xi) = \frac{1}{1 + e^{\xi/\sqrt{2D}}}.
\end{equation}

\textbf{Vérification:}
On peut vérifier (par calcul) que ce profil satisfait:
\begin{itemize}[leftmargin=1.5em]
  \item $\phi(-\infty) \to 1$ et $\phi(+\infty) \to 0$.
  \item $\phi''(\xi) = -\frac{1}{2D}\phi'(\xi)(1-2\phi(\xi))$.
  \item Lorsque substitué dans l'équation du front, il donne une approximation de très bonne qualité.
\end{itemize}

\textbf{Interprétation:} Le profil est une fonction logistique (courbe en S inversée) 
de largeur caractéristique $\sqrt{2D}$. Plus la diffusion est forte, plus le front s'étale.

\subsection*{Stabilité orbitale du front}
Un résultat central est que les fronts voyageurs sont \textbf{orbitalement stables}.
Cela signifie que si on perturbe légèrement la condition initiale autour du profil du front,
la solution reste proche de ce profil et se propage à une vitesse proche de $c^*$.
Cette stabilité est essentielle pour justifier l'observation des fronts en simulations numériques.

\section{Regularite, borne et existence de solutions globales}
Au-delà des solutions de fronts voyageurs, on s'intéresse aux propriétés générales des solutions.

\subsection*{Existence et unicité}
Dans un cadre classique (Sobolev ou Hölder), la diffusion et la continuité de $f$ permettent d'obtenir une
\textbf{existence locale} de solution faible ou classique selon la régularité de $u_0$ \cite{EvansGariepy2015}.
Pour une donnée initiale $u_0 \in L^\infty(\Omega)$, il existe $T_* > 0$ tel qu'existe une unique solution
faible en temps court $[0, T_*]$.

\subsection*{Principe du maximum et bornitude globale}
La propriété cruciale est le \textbf{principe du maximum} \cite{FifeMcLeod1977}:
Si $u_0(x) \in [0,1]$, alors la solution $u(x,t) \in [0,1]$ pour tout $t \geq 0$ et $x \in \Omega$.

\textbf{Conséquences:}
\begin{enumerate}[leftmargin=2em]
  \item La bornitude a priori permet de \textbf{prolonger la solution à tout temps} $t \in [0,\infty)$.
  \item La solution reste confinée dans une région admissible $([\![0,1]\!])$ physiquement pertinente.
  \item Cette propriété est fondamentale pour la stabilité des schémas numériques.
\end{enumerate}

\subsection*{Importance pour les simulations numériques}
Numériquement, le principe du maximum est critial:
\begin{itemize}[leftmargin=1.5em]
  \item Un schéma \textbf{qui viole le principe du maximum} produit des oscillations parasites
        et des dépassements non physiques (overshoot) hors de $[0,1]$.
  \item Les schémas \textbf{préservant le maximum} (monotones ou antidiffusifs) conservent
        les propriétés physiques même sur maillages grossiers.
  \item Cette borne assure aussi la convergence de critères d'arrêt pour les solveurs implicites.
\end{itemize}

\chapter{Partie II -- Analyse numerique}
L'analyse numerique des equations de reaction-diffusion necessite une discretisation spatiale
et temporelle \cite{Strikwerda2004,LeVeque2007}. Nous presentons plusieurs schemas explicites
et implicites.

\section{Mise en place de la discretisation}
On introduit une grille spatiale reguliere
\begin{equation}
 x_i=i\Delta x, \qquad i=0,\ldots,N,
\end{equation}
et une grille temporelle
\begin{equation}
 t^n=n\Delta t, \qquad n\ge 0.
\end{equation}
On note $u_i^n \approx u(x_i,t^n)$.

Le Laplacien est approxime par differences finies centrees:
\begin{equation}
 \Delta u(x_i,t^n) \approx \frac{u_{i+1}^n-2u_i^n+u_{i-1}^n}{(\Delta x)^2}.
\end{equation}
Sous conditions de Neumann homogenes, on peut utiliser des points fantomes:
\begin{equation}
 u_{-1}^n=u_1^n, \qquad u_{N+1}^n=u_{N-1}^n.
\end{equation}

On definit aussi
\begin{equation}
 F_i^n = D\,\Delta_h u_i^n + f(u_i^n),
\end{equation}
ou $\Delta_h$ designe le Laplacien discret.

\section{Methode 1 : Runge-Kutta d'ordre 4}
La methode de Runge-Kutta d'ordre 4 (RK4) est construite a partir d'une formule de quadrature
numerique. Elle resout l'equation differentielle ordinaire
\begin{equation}
\frac{du}{dt} = G(u,t)
\end{equation}
par une combinaison ponderee de termes evalues a differents points.

\subsection*{Derivation et formule}
Considerons le probleme initial:
\begin{equation}
\frac{du}{dt} = D\,\Delta_h u + f(u), \qquad u(t^n)=U^n.
\end{equation}
On approche l'integrale exacte
\begin{equation}
u(t^{n+1})=u(t^n) + \int_{t^n}^{t^{n+1}} G(u(t)) dt
\end{equation}
par une formule de Simpson (quadrature a 4 points). La methode RK4 utilise les
evaluations successives du second membre a 4 etapes:

\textbf{Etape 1:} evaluer $G$ a $t^n$:
\begin{equation}
 K_1 = G(U^n).
\end{equation}

\textbf{Etape 2:} evaluer $G$ en un point intermediaire $t^n + \frac{\Delta t}{2}$:
\begin{equation}
 K_2 = G\left(U^n + \frac{\Delta t}{2}K_1\right).
\end{equation}

\textbf{Etape 3:} evaluer $G$ a nouveau en $t^n + \frac{\Delta t}{2}$, mais avec une meilleure prediction:
\begin{equation}
 K_3 = G\left(U^n + \frac{\Delta t}{2}K_2\right).
\end{equation}

\textbf{Etape 4:} evaluer $G$ au point final $t^{n+1}$:
\begin{equation}
 K_4 = G\left(U^n + \Delta t\, K_3\right).
\end{equation}

La solution au temps $t^{n+1}$ est approximee par:
\begin{equation}
 U^{n+1} = U^n + \frac{\Delta t}{6}(K_1 + 2K_2 + 2K_3 + K_4).
\end{equation}

\subsection*{Application a l'equation de Nagumo discretisee}
Dans notre contexte avec discretisation spatiale, on pose $G(U)=AU+F(U)$ ou
$A$ est la matrice du Laplacien et $F$ represente la reaction nonlineaire.

Le schema RK4 devient:
\begin{align}
 K_1 &= AU^n + F(U^n), \\
 K_2 &= A\left(U^n + \frac{\Delta t}{2}K_1\right) + F\left(U^n + \frac{\Delta t}{2}K_1\right), \\
 K_3 &= A\left(U^n + \frac{\Delta t}{2}K_2\right) + F\left(U^n + \frac{\Delta t}{2}K_2\right), \\
 K_4 &= A\left(U^n + \Delta t\, K_3\right) + F\left(U^n + \Delta t\, K_3\right),
\end{align}

et la mise a jour:
\begin{equation}
 U^{n+1} = U^n + \frac{\Delta t}{6}(K_1 + 2K_2 + 2K_3 + K_4).
\end{equation}

\subsection*{Proprietes et ordre de convergence}
La methode RK4 possede les caracteristiques suivantes:

\textbf{Ordre temporel:} L'erreur locale de troncature est $O((\Delta t)^5)$, d'ou un ordre
global de convergence de \textbf{4}. C'est un progres majeur par rapport aux schemas
d'ordre 1 (Euler).

\textbf{Stabilite:} Le schema reste conditionnellement stable, mais le domaine de stabilite
(region de stabilite dans le plan complexe) est plus grand que celui d'Euler. C'est un avantage
La condition CFL pratique pour RK4 est moins restrictive que pour Euler, permettant des
pas de temps plus grands a maillage egal.

\textbf{Cout par pas:} Chaque iteration RK4 demande 4 evaluations de $G$, soit environ 4 fois
le cout d'Euler explicite par pas. Cependant, l'ordre 4 compense souvent ce surco\^{u}t en
permettant des pas plus grands et en reduisant le nombre total d'iterations.

\subsection*{Avantages et inconvenients}
\begin{itemize}[leftmargin=1.5em]
  \item \textbf{Avantages:}
    \begin{itemize}[leftmargin=1.5em]
      \item ordre 4 en temps $\Rightarrow$ bonne precision pour pas de temps moderes;
      \item stabilite meilleure qu'un schema explicite simple (Euler);
      \item methode classique et bien comprise;
      \item excellente methode de reference pour comparaisons numeriques.
    \end{itemize}
  \item \textbf{Inconvenients:}
    \begin{itemize}[leftmargin=1.5em]
      \item encore conditionnellement stable (CFL necessite);
      \item cout de 4 evaluations par pas;
      \item moins robuste qu'un schema implicite pour tres long terme.
    \end{itemize}
\end{itemize}

\section{Methode 2 : Euler implicite}
Le schema implicite est
\begin{equation}
\frac{u_i^{n+1}-u_i^n}{\Delta t} = D\frac{u_{i+1}^{n+1}-2u_i^{n+1}+u_{i-1}^{n+1}}{(\Delta x)^2} + f(u_i^{n+1}).
\end{equation}
Il necessite la resolution d'un systeme non lineaire a chaque pas.
En pratique, on peut utiliser:
\begin{itemize}[leftmargin=1.5em]
  \item Newton ou quasi-Newton;
  \item une linearisation de $f$;
  \item un traitement semi-implicite du terme non lineaire.
\end{itemize}

Son interet est sa robustesse: la partie diffusive est traitee de facon inconditionnellement
stable. C'est souvent la methode de reference quand le maillage devient fin.

\section{Methode 3 : Adams-Bashforth 2}
Le schema multipas explicite d'ordre 2 est
\begin{equation}
 u_i^{n+1} = u_i^n + \Delta t\left(\frac{3}{2}F_i^n - \frac{1}{2}F_i^{n-1}\right),
\end{equation}
ou $F_i^n = D\Delta_h u_i^n + f(u_i^n)$.

Ce schema offre un meilleur ordre temporel (2 vs 4) mais reste conditionnellement
stable. Il requiert deux niveaux temporels, d'ou une initialisation generale par Euler
ou par RK4.

\section{Methode 4 : Adams-Moulton d'ordre 2}
La methode d'Adams-Moulton est une methode multipas implicite construite par integration
de la formule d'interpolation de Lagrange. Pour l'equation
\begin{equation}
\frac{du}{dt} = G(u),
\end{equation}
on derive le schema a partir de l'integration exacte
\begin{equation}
u(t^{n+1})=u(t^n) + \int_{t^n}^{t^{n+1}} G(u(t)) dt.
\end{equation}
En utilisant une formule de quadrature (trapeze) sur cet integrale, on obtient
\begin{equation}
u^{n+1} = u^n + \frac{\Delta t}{2}\left(G(u^n)+G(u^{n+1})\right).
\end{equation}
Appliquer a notre equation de reaction-diffusion discretisee spatialement:
\begin{equation}
u_i^{n+1} = u_i^n + \frac{\Delta t}{2}\left(F_i^n + F_i^{n+1}\right),
\end{equation}
ou $F_i^n = D\Delta_h u_i^n + f(u_i^n)$.

Cette formulation peut aussi s'ecrire de maniere explicite en reorganisant:
\begin{equation}
u_i^{n+1} - \frac{\Delta t}{2}D\Delta_h u_i^{n+1} = u_i^n + \frac{\Delta t}{2}\left(D\Delta_h u_i^n + f(u_i^n) + f(u_i^{n+1})\right).
\end{equation}
Sous forme matricielle, cela devient
\begin{equation}
\left(I - \frac{\Delta t}{2}A\right)U^{n+1} = U^n + \frac{\Delta t}{2}\left(AU^n + F(U^n) + F(U^{n+1})\right),
\end{equation}
ou $A$ est la matrice du Laplacien discret.

La resolution requiert une iteration implicite (Newton ou point fixe) puisque $F(U^{n+1})$
est non lineaire.

Avantages:
\begin{itemize}[leftmargin=1.5em]
  \item ordre 2 en temps (meilleur qu'Euler implicite);
  \item stabilite inconditionnelle pour la partie diffusive;
  \item traitement equilibre entre les niveaux $n$ et $n+1$.
\end{itemize}

Cette methode est particulierement efficace pour les problemes raides ou lorsqu'on souhaite
des pas de temps plus grands sans perte de precision.

\section{Tableau recapitulatif}
\begin{table}[H]
\centering
\begin{tabular}{@{}lccccl@{}}
\toprule
Methode & Pas & Type & Ordre temps & Cout par pas & Stabilite \\
\midrule
Runge-Kutta 4 & 1 & explicite & 4 & moyen & conditionnelle \\
Euler implicite & 1 & implicite & 1 & moyen & inconditionnelle \\
Adams-Bashforth 2 & multi & explicite & 2 & faible & conditionnelle \\
Adams-Moulton 2 & multi & implicite & 2 & moyen & inconditionnelle \\
\bottomrule
\end{tabular}
\caption{Comparaison synthetique des quatre schemas consideres.}
\end{table}

\chapter{Partie III -- Etude de stabilite}
L'analyse de stabilite des schemas numeriques est cruciale pour garantir la fiabilite
des simulations \cite{Hundsdorfer1995,Rannacher1984}.

\section{Stabilite de l'equation linearisee}
On linearise autour d'un equilibre $u^*$ en posant $u=u^*+v$ avec $|v|\ll 1$.
On obtient
\begin{equation}
\partial_t v = D\partial_{xx}v + f'(u^*)v.
\end{equation}
Si l'on cherche une solution de Fourier
\begin{equation}
 v(x,t)=\rho(t)e^{ikx},
\end{equation}
alors
\begin{equation}
 \rho'(t)=\left(f'(u^*)-Dk^2\right)\rho(t).
\end{equation}
Le taux de croissance est donc
\begin{equation}
 \lambda(k)=f'(u^*)-Dk^2.
\end{equation}
On en deduit:
\begin{itemize}[leftmargin=1.5em]
  \item stabilite locale de $u^*=0$ et $u^*=1$;
  \item instabilite de $u^*=a$;
  \item attenuation des modes haute frequence par la diffusion.
\end{itemize}

\section{Analyse de Von Neumann -- Runge-Kutta d'ordre 4}
L'analyse de stabilite de Von Neumann \cite{Trefethen1996,LeVeque2007} est une technique
standard pour etudier les schemas explicites.
On considere la partie lineaire diffusee
\begin{equation}
 u_t = D u_{xx}.
\end{equation}
En posant
\begin{equation}
 u_i^n = \rho^n e^{\mathrm{i}k x_i},
\end{equation}
le schema de Runge-Kutta d'ordre 4 donne
\begin{equation}
 \rho = 1 - 4\mu\sin^2\left(\frac{k\Delta x}{2}\right),
\end{equation}
ou
\begin{equation}
 \mu = \frac{D\Delta t}{(\Delta x)^2}.
\end{equation}
La condition $|\rho|\le 1$ pour tout $k$ impose
\begin{equation}
 \mu \le \frac{1}{2},
\end{equation}
soit la condition CFL diffusive classique:
\begin{equation}
 \Delta t \le \frac{(\Delta x)^2}{2D}.
\end{equation}

\section{Analyse de Von Neumann -- Euler implicite}
Pour Euler implicite sur le meme probleme lineaire, on obtient
\begin{equation}
 \rho = \frac{1}{1+4\mu\sin^2\left(\frac{k\Delta x}{2}\right)}.
\end{equation}
Comme le denominateur est toujours strictement positif, on a automatiquement
\begin{equation}
 |\rho|\le 1 \qquad \forall \Delta t>0.
\end{equation}
Le schema est donc inconditionnellement stable pour la diffusion lineaire.

\section{Extension a AB2 et Adams-Moulton}
Pour AB2, le polynome caracteristique montre un domaine de stabilite limite. Le schema
reste donc conditionne par le maillage et par la raideur du probleme.

Pour Adams-Moulton d'ordre 2, la methode possede d'excellentes proprietes de stabilite.
Analyse de Von Neumann sur la partie diffusive (voir section suivante) montre que le schema
est A-stable et offre une error de troncature d'ordre $O((\Delta t)^3)$ pour la partie temporelle.
En pratique, c'est un schema beaucoup plus robuste qu'AB2, avec un bon compromis entre
précision et cout de calcul.

\section{Stabilite non lineaire}
Le terme de reaction influe sur le seuil de stabilite globale.
Autour d'une solution lisse, une estimation pratique est
\begin{equation}
 \Delta t \leq \min\left(\frac{(\Delta x)^2}{2D},\; \frac{1}{\max |f'(u)|}\right).
\end{equation}
Cette borne n'est pas toujours stricte mathematiquement, mais elle constitue un bon critere
pratique pour les calculs explicites.

\section{Ordre de convergence numerique}
L'ordre de convergence se mesure en comparant une solution numerique a une solution de
reference sur plusieurs maillages:
\begin{equation}
 E_h = \|u_h-u_{ref}\|.
\end{equation}
En representation log-log, on s'attend a une pente proche de 4 pour RK4,
2 pour AB2, 1 pour Euler implicite et 2 pour Adams-Moulton, tant que l'erreur spatiale ne domine pas.

\chapter{Partie IV -- Simulations et resultats}
\section{Parametres de simulation}
On peut utiliser par exemple:
\begin{equation}
 L=60, \qquad D=1, \qquad a=0.25, \qquad T=25.
\end{equation}
La condition initiale choisie est un front lisse centre en $x_0$ avec largeur $\sigma$.
Ce choix permet de visualiser clairement la propagation et d'extraire la vitesse du front.

\section{Comparaison des quatre schemas}
Les simulations doivent comparer:
\begin{itemize}[leftmargin=1.5em]
  \item le profil final $u(x,T)$;
  \item l'evolution temporelle du front $x_{0.5}(t)$;
  \item la vitesse numerique observee;
  \item la stabilite du profil et l'absence d'oscillations parasites.
\end{itemize}

RK4 offre un excellent compromis entre precision et cout. Euler implicite et Adams-Moulton sont mieux adaptes
aux pas de temps plus grands. AB2 constitue une option intermediaire avec ordre 2, moins precis que RK4
mais moins couteux et restant explicite.

\section{Etude de sensibilite}
La variation de $a$ controle le sens de propagation du front:
\begin{itemize}[leftmargin=1.5em]
  \item $a<1/2$: propagation de 1 vers 0;
  \item $a>1/2$: propagation de 0 vers 1;
  \item $a=1/2$: front quasi stationnaire.
\end{itemize}

La variation de $D$ modifie la largeur de l'interface:
\begin{itemize}[leftmargin=1.5em]
  \item grand $D$: front plus large et plus lisse;
  \item petit $D$: front plus abrupt et plus raide numeriquement.
\end{itemize}

\section{Cout computationnel}
Le cout d'un calcul depend principalement du nombre de pas de temps requis.
Les methodes explicites sont peu couteuses par pas mais demandent souvent un pas de temps
beaucoup plus petit.
Les schemas implicites sont plus chers a chaque iteration, mais peuvent reduire fortement
le nombre total d'iterations.

\section{Figures a produire}
Pour une presentation complete, on prevoit les figures suivantes:
\begin{itemize}[leftmargin=1.5em]
  \item profil initial;
  \item profils a plusieurs instants;
  \item suivi du front $u=0.5$;
  \item comparaison des vitesses numeriques et theoriques;
  \item courbes de convergence;
  \item energie discrete en fonction du temps.
\end{itemize}

\chapter{Code MATLAB}
% CORRECTION : lstinputlisting remplace par un listing inline
% car le fichier nagumo_simulation.m n'est pas disponible separement.
\begin{lstlisting}[style=matlabstyle,caption={Script MATLAB principal : nagumo\_simulation.m}]
%% nagumo_simulation.m
% Simulation de l'equation de Nagumo 1D par quatre schemas
% u_t = D*u_xx + u*(1-u)*(u-a),  x in [0,L], t in [0,T]
% Conditions de Neumann homogenes aux deux bords.

clear; close all; clc;

%% -- Parametres physiques --
L  = 60;      % longueur du domaine
D  = 1;       % coefficient de diffusion
a  = 0.25;    % parametre de seuil (0 < a < 1)
T  = 25;      % temps final

%% -- Discretisation spatiale --
N  = 200;                        % nombre d'intervalles
dx = L / N;                      % pas d'espace
x  = (0:N)' * dx;                % vecteur des noeuds (N+1 points)

%% -- Condition initiale : front lisse centre en x0 --
x0    = L / 2;
sigma = 2;
u0    = 1 ./ (1 + exp((x - x0) / sigma));

%% -- Terme de reaction --
f    = @(u) u .* (1 - u) .* (u - a);
fpri = @(u) -3*u.^2 + 2*(1+a)*u - a;   % derivee de f

%% -- Matrice Laplacien discret avec Neumann (points fantomes) --
e   = ones(N+1, 1);
A   = spdiags([e, -2*e, e], -1:1, N+1, N+1);
% Conditions de Neumann : u_{-1}=u_1 et u_{N+1}=u_{N-1}
A(1,   2)   = 2;
A(end, end-1) = 2;
A = (D / dx^2) * A;

%% -- Pas de temps : condition CFL pour les schemas explicites --
mu_max = 0.45;                        % < 0.5 pour la securite
dt_exp = mu_max * dx^2 / D;
Nt_exp = ceil(T / dt_exp);
dt_exp = T / Nt_exp;

% Pas de temps plus grand autorise pour les schemas implicites
dt_imp = 10 * dt_exp;
Nt_imp = ceil(T / dt_imp);
dt_imp = T / Nt_imp;

fprintf('dx = %.4f | dt_exp = %.6f (Nt=%d) | dt_imp = %.4f (Nt=%d)\n',...
        dx, dt_exp, Nt_exp, dt_imp, Nt_imp);

%% =========================================================
%%  SCHEMA 1 : Runge-Kutta d'ordre 4
%% =========================================================
u = u0;
front_RK4 = zeros(Nt_exp, 1);
for n = 1:Nt_exp
    % Etape 1
    K1 = A*u + f(u);
    % Etape 2
    K2 = A*(u + 0.5*dt_exp*K1) + f(u + 0.5*dt_exp*K1);
    % Etape 3
    K3 = A*(u + 0.5*dt_exp*K2) + f(u + 0.5*dt_exp*K2);
    % Etape 4
    K4 = A*(u + dt_exp*K3) + f(u + dt_exp*K3);
    % Mise a jour
    u = u + (dt_exp/6.0) * (K1 + 2*K2 + 2*K3 + K4);
    
    % Position du front (isovaleur 0.5)
    idx = find(u(1:end-1) >= 0.5 & u(2:end) < 0.5, 1, 'first');
    if ~isempty(idx)
        front_RK4(n) = x(idx) + dx * (0.5 - u(idx)) / (u(idx+1) - u(idx));
    end
end
u_RK4 = u;

%% =========================================================
%%  SCHEMA 2 : Euler implicite (semi-implicite : f explicite)
%% =========================================================
I = speye(N+1);
u = u0;
front_EI = zeros(Nt_imp, 1);
for n = 1:Nt_imp
    % (I - dt*A)*u^{n+1} = u^n + dt*f(u^n)
    rhs = u + dt_imp * f(u);
    u   = (I - dt_imp * A) \ rhs;
    idx = find(u(1:end-1) >= 0.5 & u(2:end) < 0.5, 1, 'first');
    if ~isempty(idx)
        front_EI(n) = x(idx) + dx * (0.5 - u(idx)) / (u(idx+1) - u(idx));
    end
end
u_EI = u;

%% =========================================================
%%  SCHEMA 3 : Adams-Bashforth 2
%%  Initialisation par un pas Euler implicite (pour AB2)
%% =========================================================
u     = u0;
Fold  = A*u + f(u);
u     = u + dt_exp * Fold;          % premier pas Euler
front_AB2 = zeros(Nt_exp, 1);
for n = 2:Nt_exp
    Fnew = A*u + f(u);
    u    = u + dt_exp * (1.5*Fnew - 0.5*Fold);
    Fold = Fnew;
    idx = find(u(1:end-1) >= 0.5 & u(2:end) < 0.5, 1, 'first');
    if ~isempty(idx)
        front_AB2(n) = x(idx) + dx * (0.5 - u(idx)) / (u(idx+1) - u(idx));
    end
end
u_AB2 = u;

%% =========================================================
%%  SCHEMA 4 : Adams-Moulton d'ordre 2 (semi-implicite : f explicite)
%%  Initialisation par un pas Euler implicite
%% =========================================================
u_prev = u0;
rhs0   = u_prev + dt_imp * f(u_prev);
u_cur  = (I - dt_imp * A) \ rhs0;   % premier pas Euler implicite
front_AM2 = zeros(Nt_imp, 1);
for n = 2:Nt_imp
    % (3/2*I - dt*A)*u^{n+1} = 2*u^n - 1/2*u^{n-1} + dt*f(u^n)
    rhs    = 2*u_cur - 0.5*u_prev + dt_imp * f(u_cur);
    u_new  = (1.5*I - dt_imp * A) \ rhs;
    u_prev = u_cur;
    u_cur  = u_new;
    idx = find(u_cur(1:end-1) >= 0.5 & u_cur(2:end) < 0.5, 1, 'first');
    if ~isempty(idx)
        front_AM2(n) = x(idx) + dx*(0.5 - u_cur(idx))/(u_cur(idx+1) - u_cur(idx));
    end
end
u_AM2 = u_cur;

%% =========================================================
%%  Vitesse theorique
%% =========================================================
c_th = sqrt(2*D) * (0.5 - a);
fprintf('Vitesse theorique : c = %.4f\n', c_th);

%% =========================================================
%%  Figures
%% =========================================================

% -- Profils finaux --
figure;
plot(x, u0,    'k--',  'LineWidth', 1.2, 'DisplayName', 'CI');
hold on;
plot(x, u_RK4, 'b-',   'LineWidth', 1.5, 'DisplayName', 'RK4');
plot(x, u_EI,  'r-',   'LineWidth', 1.5, 'DisplayName', 'Euler impl.');
plot(x, u_AB2, 'g-',   'LineWidth', 1.5, 'DisplayName', 'AB2');
plot(x, u_AM2,'m-',   'LineWidth', 1.5, 'DisplayName', 'Adams-Moulton');
xlabel('x'); ylabel('u(x,T)');
title('Profil final u(x,T) -- comparaison des schemas');
legend('Location', 'best'); grid on;

% -- Suivi du front --
t_exp = (1:Nt_exp)' * dt_exp;
t_imp = (1:Nt_imp)' * dt_imp;
figure;
plot(t_exp, front_RK4,  'b-',  'LineWidth', 1.5, 'DisplayName', 'RK4'); hold on;
plot(t_imp, front_EI,   'r-',  'LineWidth', 1.5, 'DisplayName', 'Euler impl.');
plot(t_exp, front_AB2,  'g.',  'DisplayName', 'AB2');
plot(t_imp, front_AM2, 'm-',  'LineWidth', 1.5, 'DisplayName', 'Adams-Moulton');
plot([0 T], [L/2 L/2 + c_th*T], 'k--', 'LineWidth', 1.2, 'DisplayName', ...
    sprintf('Theorique c=%.3f', c_th));
xlabel('t'); ylabel('x_{front}');
title('Position du front au cours du temps');
legend('Location', 'best'); grid on;
\end{lstlisting}

\chapter*{Conclusion}
\addcontentsline{toc}{chapter}{Conclusion}
L'equation de Nagumo est un excellent modele pour comprendre les fronts bistables \cite{Grindrod1988}.
L'analyse theorique montre le role du seuil $a$ et explique la vitesse de propagation.
L'etude numerique met en evidence la difference fondamentale entre les schemas explicites,
conditionnellement stables, et les schemas implicites, plus robustes.

Du point de vue pratique, RK4 offre un bon compromis entre precision et cout. Euler implicite
et Adams-Moulton sont preferables pour des simulations plus longues ou plus raides.
AB2 constitue un compromis interessant lorsque l'on veut conserver un cout faible tout
en gagnant en ordre temporel.

\chapter*{Annexes}
\addcontentsline{toc}{chapter}{Annexes}
\section*{Annexe A -- Formules utiles}
\begin{align*}
 f(u)   &= u(1-u)(u-a), \\
 f'(u)  &= -3u^2+2(1+a)u-a, \\
 c_{\mathrm{th}} &= \sqrt{2D}\left(\frac12-a\right).
\end{align*}

\section*{Annexe B -- Ordre des methodes}
\begin{longtable}{@{}p{4cm}p{2cm}p{2cm}p{2cm}p{4cm}@{}}
\toprule
Methode & Type & Ordre & Stabilite & Remarque \\
\midrule
Runge-Kutta 4 & explicite & 4 & conditionnelle & bon compromis precision/cout \\
Euler implicite & implicite & 1 & inconditionnelle & robuste pour la diffusion \\
AB2 & explicite & 2 & conditionnelle & necessite deux pas \\
Adams-Moulton & implicite & 2 & inconditionnelle & bon choix pour problemes raides \\
\bottomrule
\end{longtable}

\section*{Annexe C -- Compilation du document}
Pour compiler ce document (sans fichier \texttt{.bib} externe), utiliser simplement:
\begin{verbatim}
pdflatex nagumo_modeling.tex
pdflatex nagumo_modeling.tex
\end{verbatim}
Les references sont integrees directement dans le source.

%% =========================================================
%%  BIBLIOGRAPHIE inline (remplace nagumo.bib + natbib)
%% =========================================================
\begin{thebibliography}{99}

\bibitem{Nagumo1962}
J.~Nagumo, S.~Arimoto, S.~Yoshizawa,
\textit{An active pulse transmission line simulating nerve axon},
Proc. IRE, 50(10):2061--2070, 1962.

\bibitem{MurrayMathematicalBiology}
J.~D.~Murray,
\textit{Mathematical Biology I: An Introduction},
3rd ed., Springer, 2002.

\bibitem{Volpert2000}
A.~Volpert, V.~Volpert, V.~Volpert,
\textit{Traveling Wave Solutions of Parabolic Systems},
American Mathematical Society, 1994.

\bibitem{Aronson1978}
D.~G.~Aronson, H.~F.~Weinberger,
\textit{Multidimensional nonlinear diffusion arising in population genetics},
Adv. Math., 30(1):33--76, 1978.

\bibitem{PeletierSerafini2012}
L.~A.~Peletier, W.~C.~Troy,
\textit{Spatial Patterns: Higher Order Models in Physics and Mechanics},
Birkh\"auser, 2001.

\bibitem{EvansGariepy2015}
L.~C.~Evans,
\textit{Partial Differential Equations},
2nd ed., American Mathematical Society, 2010.

\bibitem{FifeMcLeod1977}
P.~C.~Fife, J.~B.~McLeod,
\textit{The approach of solutions of nonlinear diffusion equations to
travelling front solutions},
Arch. Rational Mech. Anal., 65(4):335--361, 1977.

\bibitem{Strikwerda2004}
J.~C.~Strikwerda,
\textit{Finite Difference Schemes and Partial Differential Equations},
2nd ed., SIAM, 2004.

\bibitem{LeVeque2007}
R.~J.~LeVeque,
\textit{Finite Difference Methods for Ordinary and Partial Differential Equations},
SIAM, 2007.

\bibitem{Ascher2008}
U.~M.~Ascher, L.~R.~Petzold,
\textit{Computer Methods for Ordinary Differential Equations and
Differential-Algebraic Equations},
SIAM, 1998.

\bibitem{Trefethen1996}
L.~N.~Trefethen,
\textit{Finite Difference and Spectral Methods for Ordinary and
Partial Differential Equations},
Cornell University, 1996.

\bibitem{Hundsdorfer1995}
W.~Hundsdorfer, J.~G.~Verwer,
\textit{Numerical Solution of Time-Dependent Advection-Diffusion-Reaction Equations},
Springer, 2003.

\bibitem{Rannacher1984}
R.~Rannacher,
\textit{Finite element solution of diffusion problems with irregular data},
Numer. Math., 43(2):309--327, 1984.

\bibitem{Grindrod1988}
P.~Grindrod,
\textit{The Theory and Applications of Reaction-Diffusion Equations},
2nd ed., Oxford University Press, 1996.

\end{thebibliography}

\end{document}